\documentclass[11pt,a4paper]{article}
% Compile with XeLaTeX in Overleaf
\usepackage[a4paper,margin=0.8in]{geometry}
\usepackage{fontspec}
\usepackage{array,tabularx,multirow}
\usepackage{amsmath,amssymb}
\usepackage{enumitem}
\setlength{\parindent}{0pt}
\setlength{\parskip}{4pt}
\renewcommand{\arraystretch}{1.45}
\setlist[itemize]{leftmargin=1.5em,itemsep=1pt,topsep=2pt}
\pagestyle{empty}
\newfontfamily\hindifont[Script=Devanagari]{Noto Serif Devanagari}

\begin{document}
\begin{center}
{\Large\bfseries KESHAV MAHAVIDYALAYA}\\[-2pt]
{\normalsize University of Delhi}
\end{center}

\vspace{8pt}
17 November 2026

\vspace{8pt}
To,\\
The Principal\\
Keshav Mahavidyalaya

\textbf{Subject: Importance of Mathematics and Technology}

Respected Sir/Madam,

Mathematics plays an important role in education, science and technology.
It develops logical thinking and problem-solving skills.

\textbf{1. Important Mathematical Concepts}
\begin{itemize}
\item Algebra and Equations
\item \underline{Geometry and Calculus}
\end{itemize}

\textbf{2. Mathematical Expressions}

The following table presents important mathematical concepts and their applications.

\begin{table}[ht]
\centering
\small
\caption{\textbf{Mathematical Expressions and Applications}}
\label{tab:math-applications}
\begin{tabular}{|>{\raggedright\arraybackslash}p{0.17\linewidth}|
                >{\centering\arraybackslash}p{0.19\linewidth}|
                >{\raggedright\arraybackslash}p{0.17\linewidth}|
                >{\centering\arraybackslash}p{0.27\linewidth}|}
\hline
\multicolumn{4}{|c|}{\textbf{Mathematical Concepts and Applications}}\\
\hline
\textbf{Algebra} & \(x^2+y^2\) & \textbf{Calculus} &
\(\displaystyle\int_0^1 x^2\,dx=\frac{1}{3}\)\\
\hline
\textbf{Geometry} & \(a^2+b^2=c^2\) & \textbf{Statistics} &
\(\mu+\sigma\)\\
\hline
\end{tabular}
\end{table}

For example, Einstein's equation is:
\[
E=mc^2 \qquad (m\in\mathbb{R})
\]

The Greek symbols \(\alpha,\beta,\theta\) are commonly used in mathematics.

Digital learning\footnote{Digital learning means the use of technology for education.}
makes mathematical concepts easier to understand.

{\hindifont हिंदी वाक्य: गणित हमारे दैनिक जीवन का महत्वपूर्ण हिस्सा है।}

I request you to encourage students to develop mathematical skills.

Thank you.

\vspace{10pt}
\begin{flushright}
Yours sincerely,\\
\textbf{Student Name}\\
Student Representative\\
Keshav Mahavidyalaya
\end{flushright}
\end{document}
